\section{企业 GenAI 典型实践}

\subsection{金融服务中的智能助理}

数字银行常见的 GenAI 路径是：把 LLM 嵌入客户服务、财富管理和内部合规流程。典型模式如下：

\begin{itemize}
    \item \textbf{输入}：用户询问投资建议 / 举报可疑交易
    \item \textbf{搜索}：实时抓取市场数据、监管规定
    \item \textbf{LLM}：生成解释性答案+操作建议
    \item \textbf{工具}：自动填报报告、生成合同草稿
\end{itemize}

\begin{knowledgebox}{金融 AI 的三层闭环}
1. 前端：LLM 生成解释性答复<br>
2. 中间层：动态搜索/检索确保信息实时<br>
3. 后端：低代码自动化工具执行命令<br>
三层之间要保证可追溯性（为什么推荐这条策略？）与风险防范（拒绝敏感操作）。
\end{knowledgebox}

\subsection{制造与供应链的 GenAI Agent}

制造公司用 GenAI Agent 连接 MES、ERP、供应链可视化工具：

\begin{itemize}
    \item 工程师上传 CAD + 用例描述
    \item Agent 读取 BOM、测试记录、故障历史，生成排查计划
    \item 直接触发工单或通知采购规划
\end{itemize}

\begin{warningbox}{测试场景必须真实再现}
很多 Agent 在测试环境表现良好，但一进真实车间就崩溃是因为数据分布差异（日志 vs. 实时 sensor）。必须在真实操作中收集 feedback 并回放给模型。
\end{warningbox}

\subsection{本章小结}

不同领域的实践都遵循同样的闭环模式：AI 提供理解+生成，搜索/数据确保事实，工具/流程完成执行，且整个链路需要可追溯和可控。

